\documentclass[12pt]{article}
\usepackage[utf8]{inputenc}
\usepackage[T1]{fontenc}
\usepackage{hyperref}
\title{Turks and Caicos English Creole Language Resources for K-12 Curriculum Development}
\author{Compiled by Claude Sonnet 5 (claude-sonnet-5) for Andr\'as Kornai}
\date{}
\begin{document}
\maketitle

\section{Introduction}
Turks and Caicos Creole (ISO 639-3 \texttt{tch}) is an English-based creole. Wikipedia describes it as a dialect of Bahamian Creole, closely related to Bermudian Creole, and not thoroughly studied. Speaker figures conflict, from about 10,700 (1995) to 34,000 (2019). It has no official status and appears to be shifting toward Caribbean English. This survey found very little material, and nothing designed for schools.

\section{Primary Educational Resources}
No primers, textbooks or school programmes for the creole turned up in the searches.

\section{Linguistic and Reference Materials}
\begin{itemize}
\item Wikipedia, Turks and Caicos Creole: \url{https://en.wikipedia.org/wiki/Turks_and_Caicos_Creole}
\item Glottolog, languages of the Turks and Caicos Islands (catalogue listing): \url{https://glottolog.org/glottolog?country=TC}
\end{itemize}
The search summary described Cecilia Cutler's chapter on English in the Turks and Caicos Islands (Grand Turk) as the first linguistic study of the islands' speech. I did not see a direct URL for it, so it is not cited as a link here.

\section{Religious and Community Resources}
None found.

\section{Digital and Audio Resources}
\begin{itemize}
\item Global Recordings Network, language entry for \texttt{tch}. It currently lists no recordings: \url{https://globalrecordings.net/en/language/tch}
\end{itemize}

\section{Curriculum Development Context}
Little published material exists. Because the language is close to Bahamian Creole, materials for Bahamian and wider Caribbean creoles may be a useful basis, but none were verified in these searches. Local documentation would be needed first.
\end{document}